\chapter{Natural Language Processing and Text Representation}
\label{chap:nlp-representation-en}

\FloatBarrier
\section{Graded Relevance Metrics (CG, DCG, NDCG)}
\label{sec:recap-ndcg-en}

In real-world Information Retrieval systems, the assumption of purely binary relevance (relevant vs. non-relevant) is inadequate to capture the nuances of search results. Relevance judgments are therefore expressed on graded scales (e.g., a Likert scale from $0$ to $3$ or $0$ to $4$, indicating increasing degrees of pertinence).

The evaluation of a ranked list of documents up to position $p$ is founded on two core principles:
\begin{enumerate}
    \item \textbf{Gain:} The higher a document with high relevance appears within the SERP (\textit{Search Engine Results Page}), the greater the utility perceived by the user.
    \item \textbf{Discounting:} Highly relevant documents placed at lower ranks deliver progressively less utility than if placed at the top, since user attention monotonically decays down the list.
\end{enumerate}

\subsection{Cumulative Gain (CG) and Discounted Cumulative Gain (DCG)}

Given a vector of graded relevance scores $\mathbf{r} = [r_1, r_2, \dots, r_p]$ for the top-$p$ documents returned for a query $q$:
\begin{definition}[Cumulative Gain at Rank $p$]
\textbf{Cumulative Gain} ($CG_p$) is the unweighted sum of relevance values up to rank $p$:
\begin{equation}
CG_p = \sum_{i=1}^{p} r_i.
\end{equation}
\textit{Inherent limitation:} $CG_p$ is position-insensitive: placing a highly relevant document ($r_i = 3$) at rank 1 or rank 10 produces identical $CG$ scores.
\end{definition}

To model user attention decay as rank $i$ increases, a logarithmic discount function is applied:

\begin{definition}[Discounted Cumulative Gain at Rank $p$]
In the standard formulation and the exponential gain variant widely adopted in industrial search engines:
\begin{align}
\text{DCG}_p &= r_1 + \sum_{i=2}^{p} \frac{r_i}{\log_2(i)}, \label{eq:dcg-std-en-vsm} \\
\text{DCG}_p^{\text{exp}} &= \sum_{i=1}^{p} \frac{2^{r_i} - 1}{\log_2(i + 1)}. \label{eq:dcg-exp-en-vsm}
\end{align}
The exponential variant~\eqref{eq:dcg-exp-en-vsm} heavily penalizes the failure to place top-relevance documents at the very highest ranks.
\end{definition}

\subsection{Normalized DCG (NDCG)}

Raw DCG has a critical structural drawback: it is strictly \textbf{query-dependent}. 
If a specific query $q_A$ has $100$ highly relevant documents ($r=3$) in the collection, its theoretical DCG will naturally be huge. If another query $q_B$ has only one relevant document in the entire corpus, its maximum attainable DCG will be very small. Consequently, raw DCG scores cannot be directly compared or averaged across different queries.

To enable cross-query comparison and global macro-averaging across an evaluation benchmark, we compute the \textbf{Ideal DCG} ($\text{IDCG}_p$), which is the DCG of the ideal ranking obtained by sorting all documents in descending order of relevance:
\begin{equation}
\text{NDCG}_p = \frac{\text{DCG}_p}{\text{IDCG}_p} \quad \in [0, 1].
\end{equation}
With normalization, $\text{NDCG}_p = 1.0$ guarantees that the search engine achieved the best possible ordering up to depth $p$.

\FloatBarrier
\section{Characteristics of Natural Language and Zipf's Law}
\label{sec:language-zipf-en}

\subsection{Properties and Inherent Complexity of Natural Language}

Processing natural language computationally presents unique hurdles due to its structural characteristics:
\begin{itemize}
    \item \textbf{Vast vocabularies:} Real-world vocabularies contain hundreds of thousands of distinct words.
    \item \textbf{Complex, hierarchical syntax:} Sentences are not mere flat sequences, but nested syntax trees with long-range grammatical dependencies.
    \item \textbf{Compositional semantics:} The overall meaning of an utterance is a function of the meanings of its parts and their syntactic arrangement.
    \item \textbf{Pervasive ambiguity:}
    \begin{itemize}
        \item \textit{Part-of-Speech (POS) ambiguity:} A single word form functions as noun or verb depending on context (e.g., \textit{beat}: \textit{``Listen to this cool beat''} [noun] vs. \textit{``I'll beat you at checkers!''} [verb]).
        \item \textit{Word sense ambiguity (Polysemy and Homonymy):} Multiple distinct meanings (e.g., \textit{interest}: \textit{``There is interest on this topic''} [curiosity] vs. \textit{``The bank raised the interest rate''} [financial charge]).
        \item \textit{Syntactic ambiguity:} Sentences with multiple valid parse trees (e.g., \textit{``I shot an elephant in my pajamas''}).
    \end{itemize}
    \item \textbf{Heavy reliance on Shared World Knowledge:}
    \begin{itemize}
        \item \textit{``I'm faster than Lewis Hamilton''}: Requires knowing Lewis Hamilton is a Formula 1 driver to interpret what speed is being compared.
        \item \textit{``Alexis and Kate are sisters''}: Asserts a symmetric kinship relation between two entities.
        \item \textit{``Alexis and Kate are mothers''}: Attributes a shared ontological property to both individuals, without implying a direct relationship between them.
        \item \textit{``My mother is younger than me''}: Syntactically well-formed, but semantically and biologically impossible.
    \end{itemize}
    \item \textbf{Temporal evolution:} Vocabularies evolve rapidly through technological innovation and socio-cultural shifts (e.g., \textit{``to google''}, \textit{``crowdfunded''}, \textit{``in the cloud''}, \textit{``super spreader''}, \textit{``wear a mask''}).
\end{itemize}

\subsection{Zipf's Law and the Principle of Least Effort}

In any natural language corpus, the empirical frequency $f$ of a term is inversely proportional to its rank $r$ when sorted by frequency:
\begin{equation}
f \propto \frac{1}{r} \implies f \cdot r \approx C,
\end{equation}
where $C$ is a corpus-specific constant.
In Mandelbrot's generalized analytical formulation (\textit{Zipf-Mandelbrot law}):
\begin{equation}
f(r) \propto \frac{1}{(r + \beta)^\alpha},
\end{equation}
typically with $\alpha \approx 1$ and $\beta \approx 2.7$.

\begin{noteblock}{Zipf's Principle of Least Effort}
George Kingsley Zipf explained this hyperbolic distribution as an equilibrium between two opposing cognitive forces:
\begin{enumerate}
    \item \textbf{Speaker's effort:} Aims to minimize cognitive strain by using a small vocabulary of frequent, short, polysemous words.
    \item \textbf{Hearer's effort:} Prefers a large vocabulary of distinct, rare, and precise words to minimize decoding ambiguity.
\end{enumerate}
The evolutionary compromise between speaker economy and listener clarity yields Zipf's law.
\end{noteblock}

From Zipf's formulation follow two key linguistic properties:
\begin{itemize}
    \item The number of meanings $m$ of a word scales inversely with the square root of its frequency: $m \propto 1 / \sqrt{f}$.
    \item An inverse relationship connects word frequency and word length: high-frequency words are almost always the shortest.
\end{itemize}

\subsection{Rare Features and Stop Words}

Zipf's distribution exposes two critical extremes for search system design:
\begin{enumerate}
    \item \textbf{Rare Features and \textit{Hapax Legomena}:} Words appearing only once in the entire corpus (\textit{hapax legomena}) typically comprise $40\%$ to $50\%$ of vocabulary $|V|$. A feature may be rare for two distinct reasons:
    \begin{itemize}
        \item It is a highly specialized technical term with \textbf{enormous discriminative content} (e.g., \textit{``arachnocentric''}).
        \item It is a spelling typo or an arbitrary URL/file slug.
    \end{itemize}
    Pruning rare terms drastically shrinks index size, but risks eliminating highly specific search needles.
    \item \textbf{Stop Words:} Extremely high-frequency syntactic glue words (\textit{``the'', ``of'', ``a'', ``to'', ``in''}). Historically stripped during indexing to save memory and boost traversal speed (\textit{``the president of the united states of america''} $\to$ \textit{``president united states america''}).
\end{enumerate}

\begin{alertblock}{Warning on Stop Word Lists}
Stop word lists are domain-specific. For instance, MySQL's default stop word list originally included words like \textit{``appreciate''}, \textit{``serious''}, and \textit{``unfortunately''}, which are vital features for Sentiment Analysis and opinion-oriented search tasks.
\end{alertblock}

\FloatBarrier
\section{Data Acquisition and Web Parsing: urllib and BeautifulSoup}
\label{sec:crawling-parsing-en}

The initial step in search engine construction is crawling and HTML text extraction, depicted in Figure~\ref{fig:web-crawling-pipeline-en}.

\begin{figure}[H]
\centering
\resizebox{0.92\linewidth}{!}{%
\begin{tikzpicture}[
    node distance=1.5cm,
    stage/.style={
        rectangle,
        rounded corners=6pt,
        thick,
        align=center,
        inner sep=8pt,
        minimum width=11.0cm,
        minimum height=1.3cm
    },
    server/.style={stage, draw=blue!70!black, fill=blue!5},
    html/.style={stage, draw=orange!80!black, fill=orange!6},
    clean/.style={stage, draw=teal!70!black, fill=teal!6},
    arr/.style={
        -{Stealth[length=3.5mm, width=2.5mm]},
        very thick,
        draw=gray!75!black
    },
    toolbadge/.style={
        midway,
        fill=white,
        inner sep=4pt,
        thick,
        rounded corners=4pt,
        font=\small
    }
]
    % Level 1: Remote Web Server
    \node[server] (server) {
        \textbf{\large\color{blue!85!black}1. Remote Web Server (Source)}\\[2pt]
        {\footnotesize HTTP/HTTPS Endpoint $\;\cdot\;$ Remote Web Resource $\;\cdot\;$ Network Connection}
    };

    % Level 2: Raw HTML
    \node[html, below=1.5cm of server] (html) {
        \textbf{\large\color{orange!85!black}2. Raw HTML Document}\\[2pt]
        {\footnotesize Full DOM tree with markup tags, XML metadata, scripts, and CSS stylesheets}
    };

    % Level 3: Clean Text
    \node[clean, below=1.5cm of html] (clean) {
        \textbf{\large\color{teal!85!black}3. Extracted Clean Text}\\[2pt]
        {\footnotesize Textual payload only $\;\cdot\;$ Stripped of markup and code $\;\cdot\;$ Ready for tokenization}
    };

    % Vertical arrows with centered tool badges
    \draw[arr] (server) -- (html)
        node[toolbadge, draw=blue!50!black] {
            \textbf{\texttt{urllib.request}} \quad {\footnotesize\color{black!75}(HTTP download and response retrieval)}
        };

    \draw[arr] (html) -- (clean)
        node[toolbadge, draw=teal!60!black] {
            \textbf{\texttt{BeautifulSoup}} (\texttt{bs4}) \quad {\footnotesize\color{black!75}(DOM tree parsing and text stripping)}
        };
\end{tikzpicture}%
}
\caption{HTTP web crawling and HTML stripping pipeline.}
\label{fig:web-crawling-pipeline-en}
\end{figure}

\subsection{Acquisition with urllib}
In Python, low-level HTTP requests are handled cleanly via \texttt{urllib.request}:

\begin{lstlisting}[style=pythoncode,caption={Fetching raw HTML payload over HTTP.}]
from urllib import request

url = "http://hpc.isti.cnr.it/~nardini/"
response = request.urlopen(url)
html = response.read().decode('utf8')
# Yields the full HTML document string with markup
\end{lstlisting}

\subsection{DOM Parsing and Stripping with BeautifulSoup}
Web pages embed scripts and formatting markup that would pollute term statistics. \texttt{BeautifulSoup} navigates the DOM tree to extract pure content:

\begin{lstlisting}[style=pythoncode,caption={HTML markup stripping with BeautifulSoup.}]
from bs4 import BeautifulSoup

soup = BeautifulSoup(html, "html5lib")

# Direct extraction of all concatenated text:
text_all = soup.get_text()

# Granular extraction targeting only paragraph (<p>) tags:
text_paragraphs = [''.join(s.findAll(text=True)) for s in soup.findAll('p')]
\end{lstlisting}

\FloatBarrier
\section{Text Preprocessing Pipeline: Tokenization and Sentence Splitting}
\label{sec:tokeniz-splitting-en}

\begin{noteblock}{Data Engineering Rule of Thumb}
In real-world IR and NLP architectures, \textbf{80\% of engineering time and effort} is spent on data ingestion, cleaning, normalization, and preprocessing pipelines before any model or indexing algorithm can be run.
\end{noteblock}

\subsection{Tokenization}

Tokenization breaks raw character streams into discrete linguistic units called \textbf{tokens}:

\begin{lstlisting}[style=pythoncode,caption={Tokenization and stop word removal using NLTK.}]
import nltk
from nltk import word_tokenize
from nltk.corpus import stopwords

text = "the president of the united states of america"
tokens = word_tokenize(text)
# Output: ['the', 'president', 'of', 'the', 'united', 'states', 'of', 'america']

stop_words = set(stopwords.words('english'))
features = set(tokens)
filtered_features = features.difference(stop_words)
# Output: {'america', 'president', 'states', 'united'}
\end{lstlisting}

\subsection{Sentence Splitting and Punctuation Ambiguity}

Segmenting unpunctuated or continuous text into standalone sentences (\textit{sentence boundary disambiguation}) is challenging due to the overloaded role of the period (\texttt{.}):
\begin{enumerate}
    \item Sentence terminator.
    \item Acronym, abbreviation, or honorific marker (\textit{U.S.A.}, \textit{e.g.}, \textit{St. Louis}, \textit{Dr.}).
    \item Numerical decimal or time separator (\textit{8.30}).
    \item Discourse ellipsis (\textit{...}).
\end{enumerate}

\begin{example}[Sentence Boundary Disambiguation]
Consider the classroom example:
\begin{center}
\textit{``I'll take the 8.30 train to St. Louis... or maybe I'll stay in LA. Time will tell!''}
\end{center}
A naive \texttt{split('.')} rule produces severe errors:
\begin{itemize}
    \item The dot in \textit{8.30} is part of a time token.
    \item The dot in \textit{St.} marks an abbreviation for \textit{Saint}.
    \item The ellipsis (\textit{...}) indicates a rhetorical pause without closing the logical clause.
    \item Only the dot after \textit{LA.} and the exclamation point in \textit{Time will tell!} are true sentence boundaries.
\end{itemize}
NLTK's \texttt{sent\_tokenize} employs the unsupervised \textit{Punkt} tokenizer, trained to identify collocations and abbreviations to preserve genuine sentence boundaries.
\end{example}

\FloatBarrier
\section{Vocabulary Construction and Counting}
\label{sec:vocabolario-conteggio-en}

\subsection{Formal Vocabulary Definition}

\begin{definition}[Collection Vocabulary]
Given a document collection $\mathcal{D} = \{d_1, d_2, \dots, d_N\}$, the \textbf{Vocabulary} (or feature set, denoted by $V$ or $F$) is the union of all distinct term types across all documents:
\begin{equation}
V = \bigcup_{d \in \mathcal{D}} \{ t \in d \}.
\end{equation}
Its cardinality $|V|$ defines the dimensionality of the vector feature space. It is strictly \textbf{dataset-dependent}: adding or removing documents alters $V$ and its dimension.
\end{definition}

\subsection{Algorithmic Analysis: Choosing the Counting Data Structure}

When designing term frequency counters, beginners often suggest using a set (\texttt{set}):
\begin{itemize}
    \item \textbf{Why a \texttt{set} fails:} A set enforces element uniqueness and discards duplicates upon first encounter, making it impossible to accumulate occurrence counts.
    \item \textbf{Correct approach:} A hash table (in Python, \texttt{dict} or \texttt{Counter}), mapping each unique token string to its running integer count.
\end{itemize}

Algorithm~\ref{alg:costruzione-vocabolario-en} details this procedure.

\begin{algorithm}[H]
\caption{Vocabulary Construction and Collection Term Frequency Counting}
\label{alg:costruzione-vocabolario-en}
\KwInput{Document collection $\mathcal{D} = \{d_1, d_2, \dots, d_N\}$}
\KwOutput{Vocabulary $V$ and frequency map $\text{freq\_dict}$}
$\text{freq\_dict} \gets \text{empty hash map}$\;
\ForEach{$d \in \mathcal{D}$}{
    $\text{tokens} \gets \text{WordTokenize}(d)$\;
    \ForEach{$t \in \text{tokens}$}{
        \eIf{$t \in \text{freq\_dict}$}{
            $\text{freq\_dict}[t] \gets \text{freq\_dict}[t] + 1$\;
        }{
            $\text{freq\_dict}[t] \gets 1$\;
        }
    }
}
$V \gets \text{keys}(\text{freq\_dict})$\;
\Return{$V, \text{freq\_dict}$}\;
\end{algorithm}

At completion, $|V| = \text{len}(\text{freq\_dict})$, and the collection frequency $cf_t$ is recorded for every term $t \in V$.

\FloatBarrier
\section{Morphological Normalization: Stemming vs. Lemmatization}
\label{sec:stemming-lemmatizzazione-en}

To collapse inflected variants into a single canonical index form, two competing paradigms exist (summarized in Table~\ref{tab:stem-vs-lemma-en}).

\begin{table}[H]
\centering
\small
\begin{tabularx}{\linewidth}{l X X}
\toprule
\textbf{Criterion} & \textbf{Stemming} & \textbf{Lemmatization} \\
\midrule
\textbf{Goal} & Strip affixes to reach a crude root (\textit{stem}). & Map word form to legitimate dictionary base (\textit{lemma}). \\
\addlinespace
\textbf{Methodology} & Rule-based heuristic suffix chopping. & Complete morphological and syntactic analysis via POS and lexical graphs. \\
\addlinespace
\textbf{Lexical Validity} & Low: often generates non-words or invalid stems. & High: always produces genuine dictionary words. \\
\addlinespace
\textbf{POS Dependency} & None: operates purely on character surface strings. & Crucial: requires contextual Part-of-Speech tagging. \\
\addlinespace
\textbf{Computational Cost} & Negligible: instantaneous string rewrites in $\BigO{1}$. & Heavy: requires statistical sequence models or graph lookups. \\
\bottomrule
\end{tabularx}
\caption{Systematic comparison between heuristic Stemming and Lemmatization.}
\label{tab:stem-vs-lemma-en}
\end{table}

\subsection{Stemming and the Porter Stemmer Failure Case}
The Porter Stemmer applies sequential suffix-stripping rules (e.g., \texttt{-ing}, \texttt{-s}, \texttt{-ed}).

\begin{lstlisting}[style=pythoncode,caption={Stemming with PorterStemmer illustrating lexical breakdown.}]
from nltk.stem.porter import PorterStemmer

stemmer = PorterStemmer()
stemmer.stem('cars')  # Result: 'car' (valid)
stemmer.stem('was')   # Result: 'wa'  (LEXICAL ERROR)
\end{lstlisting}

\begin{alertblock}{Root Cause of the Stemmer Error on ``was''}
The stemmer identifies the trailing \texttt{'s'} in \textit{``was''} and mechanically applies the English plural reduction rule. Lacking syntax awareness (that \textit{``was''} is an inflected past tense form of \textit{to be}), it produces the nonexistent pseudo-stem \texttt{'wa'}.
\end{alertblock}

\subsection{Lemmatization with WordNet}
Lemmatization queries a lexical ontology (WordNet). Without an explicit POS tag, the default assumption is that the token is a noun:

\begin{lstlisting}[style=pythoncode,caption={Context-aware lemmatization with WordNetLemmatizer.}]
from nltk.stem.wordnet import WordNetLemmatizer

lmtzr = WordNetLemmatizer()
lmtzr.lemmatize('cars')                  # Output: 'car'
lmtzr.lemmatize('was')                   # Output: 'wa' (default: noun)
lmtzr.lemmatize('was', pos='v')          # Output: 'be' (CORRECT: verb 'to be')
\end{lstlisting}

\subsection{Engineering Trade-offs in Large-Scale Search Engines}
Historically, commercial search engines processing billions of web documents have heavily favored \textbf{Stemming}: its near-zero computational overhead outweighs occasional morphological errors when balancing indexing throughput across planetary datasets.

\FloatBarrier
\section{The Bag of Words (BoW) Model}
\label{sec:bag-of-words-en}

\subsection{Transition from Sequence to Set}
Unstructured text originates as an ordered token sequence:
\begin{equation}
\text{text} \longrightarrow [t_1, t_2, t_3, \dots, t_k].
\end{equation}
The \textbf{Bag of Words (BoW)} model abstracts the text by collapsing this sequence into an unordered set (or multiset with counts):
\begin{lstlisting}[style=pythoncode,caption={Bag of Words set representation.}]
text = 'the president of the united states of america'
tokens = word_tokenize(text)      # Length 8 list
bow = set(tokens)                 # 6 unique items
# bow: {'america', 'of', 'president', 'states', 'the', 'united'}
\end{lstlisting}

\subsection{Structural Flaws of BoW}
\begin{enumerate}
    \item \textbf{Loss of local term frequencies} in pure set formulations.
    \item \textbf{Total loss of word order and syntax:} BoW assumes word placement is irrelevant to semantic content.
\end{enumerate}

\begin{example}[Semantic Collapse in Bag of Words]
Consider two sentences with diametrically opposed meanings:
\begin{itemize}
    \item $t_1$: \textit{``I won, and thus you lose.''}
    \item $t_2$: \textit{``I lose, and thus you won.''}
\end{itemize}
Evaluating their BoW representations in Python:
\begin{lstlisting}[style=pythoncode,caption={Demonstration of BoW semantic indistinguishability.}]
t1 = set(word_tokenize('I won, and thus you lose.'))
t2 = set(word_tokenize('I lose, and thus you won.'))

print(t1 == t2)  # Output: True
\end{lstlisting}
Under BoW assumptions, the two texts are identical. This demonstrates that human language semantics is deeply \textbf{compositional} and sensitive to syntax order.
\end{example}

\FloatBarrier
\section{N-grams: Word n-grams and Character n-grams}
\label{sec:ngrams-en}

To recover local context without building complete parse trees, the vocabulary is augmented with \textbf{n-grams}.

\subsection{Word n-grams}
A \textbf{word n-gram} is a contiguous sequence of $n$ words, capturing local phrase order:

\begin{lstlisting}[style=pythoncode,caption={Local syntax preservation using word bigrams.}]
import nltk

t1 = set(nltk.ngrams(nltk.word_tokenize('I won, and thus, you lose.'), 2))
t2 = set(nltk.ngrams(nltk.word_tokenize('I lose, and thus, you won.'), 2))

# t1 contains: ('won', 'and'), ('you', 'lose'), ...
# t2 contains: ('lose', 'and'), ('you', 'won'), ...
print(t1 == t2)  # Output: False
\end{lstlisting}
Bigrams ($n=2$) preserve local subject-verb combinations, successfully telling $t_1$ and $t_2$ apart.

\subsection{Character n-grams}
\textbf{Character n-grams} decompose words into contiguous character slices of length $n$ (commonly trigrams, $n=3$). This provides resilience against typos and captures morphological affinity:

\begin{lstlisting}[style=pythoncode,caption={Typo tolerance using character trigrams.}]
t1 = set(nltk.ngrams('rainbow', 3))
t2 = set(nltk.ngrams('rainbaw', 3))  # Contains typo: 'a' instead of 'o'

shared = t1.intersection(t2)
# Output: {('r', 'a', 'i'), ('a', 'i', 'n'), ('i', 'n', 'b')}
\end{lstlisting}
Despite the misspelling, the two strings share 3 out of 5 trigrams ($60\%$), allowing the search engine to detect high geometric proximity and trigger query auto-correction.

\FloatBarrier
\section{The Vector Space Model (VSM)}
\label{sec:vsm-intro-en}

In the \textbf{Vector Space Model} (VSM), documents and queries are embedded into a real vector space whose dimension matches vocabulary size:
\begin{equation}
\text{Dimensionality} = |F| = |V|.
\end{equation}
Every orthogonal axis corresponds to a distinct term $f \in V$. Documents and queries are vectors:
\begin{equation}
\vec{v}(d) \in \mathbb{R}^{|V|}, \quad \vec{v}(q) \in \mathbb{R}^{|V|}.
\end{equation}

\begin{figure}[H]
\centering
\resizebox{0.82\linewidth}{!}{%
\begin{tikzpicture}[scale=1.1, >=Stealth]
    \draw[->, thick, gray!70!black] (0,0) -- (6.0, 0) node[below] {\small Dimension 1 ($f_1$: \texttt{"weekend"})};
    \draw[->, thick, gray!70!black] (0,0) -- (0, 4.5) node[left] {\small Dimension 2 ($f_2$: \texttt{"weather"})};

    \draw[->, very thick, blue!80!black] (0,0) -- (2.0, 3.8) node[above right] {\small $\vec{v}(d_1)$ (\textit{``best weather on the weekend''})};
    \draw[->, very thick, teal!80!black] (0,0) -- (4.2, 2.2) node[right] {\small $\vec{v}(d_2)$ (\textit{``nice weather''})};

    \draw[red!80!black, thick] (0.8, 1.52) arc[start angle=62, end angle=28, radius=1.7];
    \node[red!80!black] at (1.6, 1.25) {\small $\theta$};
\end{tikzpicture}%
}
\caption{Two-dimensional geometric document vectors in the Vector Space Model.}
\label{fig:vsm-geometric-2d-en}
\end{figure}

\subsection{Formal Mapping via Canonical Unit Vectors}
Each feature $f_i \in V$ corresponds to a standard basis vector:
\begin{equation}
\vec{v}(d) = \sum_{f \in d} w_{f,d} \, \vec{v}(f),
\end{equation}
where $w_{f,d} \in \mathbb{R}$ is the weight of term $f$ in document $d$.

\subsection{Extreme Vector Sparsity}
For a typical web index with $|V| = 10^6$ and an average document length of $200$ distinct words:
\begin{equation}
\frac{|\{ i \mid v_i(d) \ne 0 \}|}{|V|} = \frac{200}{1\,000\,000} = 0.02\%.
\end{equation}
Over \textbf{99.98\%} of elements are zeros. Storing them as dense 32-bit floating-point arrays ($4\,\text{MB}$ per document) would crash system memory, necessitating specialized sparse vector formats and inverted indices.

